\documentclass[11pt]{article}
\usepackage[utf8]{inputenc}
\usepackage{geometry}
\geometry{a4paper, margin=1in}
\usepackage{hyperref}
\usepackage{graphicx}
\usepackage{tabularx}
\usepackage{booktabs}
\usepackage{enumitem}
\usepackage{amsmath}
\usepackage{float}
\usepackage{xcolor}

\definecolor{bestgreen}{RGB}{200,240,200}

\title{Landslide Identification Using ML \\ \large Lifecycle 1 Report: Foundation \& Baseline Models}
\author{M.Tech Project}
\date{2026-03-18}

\begin{document}

\maketitle

% ============================================================================
\section{Objective}
% ============================================================================

The goal of Lifecycle 1 is to establish the foundational pipeline for automated landslide identification from satellite imagery. This involves:

\begin{enumerate}
    \item Designing a dual-phase architecture: Phase 1 (CNN-based image classification) + Phase 2 (HMM-based temporal risk forecasting).
    \item Establishing a performance baseline using a simple CNN (AlexNet).
    \item Investigating whether transfer learning and modern architectures (EfficientNet-B3) can improve detection accuracy on a small satellite imagery dataset.
    \item Building an end-to-end pipeline from raw satellite image to landslide risk report.
\end{enumerate}

% ============================================================================
\section{Dataset Description}
% ============================================================================

\subsection{Phase 1: HR-GLDD (High Resolution Global Landslide Detection Dataset)}

\begin{table}[H]
    \centering
    \renewcommand{\arraystretch}{1.2}
    \begin{tabular}{lcccc}
        \toprule
        \textbf{Split} & \textbf{Landslide} & \textbf{Non-Landslide} & \textbf{Total} & \textbf{Usage} \\
        \midrule
        Train      & 616   & 1,574 & 2,190 & Model training \\
        Validation & 158   & 392   & 550   & Hyperparameter tuning \\
        Test       & 211   & 488   & 699   & Final evaluation \\
        \midrule
        \textbf{Total} & \textbf{985} & \textbf{2,454} & \textbf{3,439} & \\
        \bottomrule
    \end{tabular}
    \caption{HR-GLDD dataset splits. Source: Zenodo 7189381. Images are optical satellite imagery converted from 4-band numpy arrays to 3-band RGB JPEG.}
\end{table}

\noindent \textbf{Key challenge:} The dataset is imbalanced --- 72\% non-landslide vs 28\% landslide. The training set contains only 2,190 images, which is small for deep learning.

\subsection{Phase 2: NASA Global Landslide Catalog (GLC)}

\begin{itemize}[noitemsep]
    \item \textbf{Records:} 9,471 landslide events (1988--2016)
    \item \textbf{Coverage:} 141 countries
    \item \textbf{Sequences:} 112 country-level temporal sequences, 9,442 total observations
    \item \textbf{Features:} Landslide type, trigger mechanism, location, date
\end{itemize}

% ============================================================================
\section{Methodology}
% ============================================================================

\subsection{Dual-Phase Architecture}

The pipeline consists of two independent phases:

\begin{enumerate}
    \item \textbf{Phase 1 --- CNN Image Classification:} A convolutional neural network classifies each satellite image as \textit{landslide} or \textit{non-landslide} and outputs a confidence score.
    \item \textbf{Phase 2 --- HMM Temporal Forecasting:} If Phase 1 detects a landslide, the geographic context (country) is passed to a Hidden Markov Model trained on the NASA GLC. The HMM predicts the landslide type, occurrence probability, peak risk month, and future risk forecast.
\end{enumerate}

\noindent This separation allows each component to be improved independently in subsequent lifecycles.

\subsection{Result 0 (r0): AlexNet Trained from Scratch}

\begin{itemize}[noitemsep]
    \item \textbf{Architecture:} AlexNet --- 5 convolutional layers + 3 fully connected layers
    \item \textbf{Input:} 227$\times$227 RGB images
    \item \textbf{Training:} 30 epochs, batch size 64, LR = 0.0001, StepLR scheduler (step=10, gamma=0.1)
    \item \textbf{Parameters:} $\sim$62M (all trained from scratch)
    \item \textbf{Best epoch:} 13 (val\_acc = 78.55\%)
    \item \textbf{HMM:} 4 hidden states, 7 observation symbols (landslide type only)
    \item \textbf{Rationale:} AlexNet is a simple, well-understood architecture that serves as a performance floor. Any subsequent changes must demonstrably beat this baseline to justify added complexity.
\end{itemize}

\subsection{Result 1 (r1): AlexNet with Transfer Learning}

\begin{itemize}[noitemsep]
    \item \textbf{Architecture:} AlexNet pretrained on ImageNet (1.2M images)
    \item \textbf{Modification:} Final classifier layer replaced for 2-class output. Feature layers frozen initially, then unfrozen for fine-tuning.
    \item \textbf{Rationale:} With only 2,190 training images, training from scratch forces the model to learn basic visual features (edges, textures) from very limited data. Transfer learning provides a rich feature foundation from ImageNet, allowing the model to focus on learning landslide-specific patterns.
\end{itemize}

\subsection{Result 2 (r2): EfficientNet-B3 with Transfer Learning}

\begin{itemize}[noitemsep]
    \item \textbf{Architecture:} EfficientNet-B3 --- compound-scaled CNN with 26+ convolutional blocks
    \item \textbf{Input:} 300$\times$300 RGB images (native resolution, 74\% more pixels than AlexNet's 227$\times$227)
    \item \textbf{Training:} Feature layers frozen for first 5 epochs (warm-up), then fully unfrozen with differential learning rates (head: 1e-4, features: 1e-5). CosineAnnealingLR scheduler.
    \item \textbf{Parameters:} $\sim$12M (5$\times$ fewer than AlexNet despite being far deeper)
    \item \textbf{Best epoch:} 24 (val\_acc = 79.82\%)
    \item \textbf{HMM:} Upgraded to 6 hidden states
    \item \textbf{Rationale:} AlexNet's shallow architecture (5 conv layers) cannot capture fine-grained spatial textures. EfficientNet-B3 uses compound scaling (depth $\times$ width $\times$ resolution) for superior feature extraction with fewer parameters.
\end{itemize}

% ============================================================================
\section{Results}
% ============================================================================

\subsection{Metrics Comparison}

\begin{table}[H]
    \centering
    \renewcommand{\arraystretch}{1.2}
    \begin{tabular}{llccccc}
        \toprule
        \textbf{Result} & \textbf{Architecture} & \textbf{Acc (\%)} & \textbf{Prec (\%)} & \textbf{Rec (\%)} & \textbf{F1 (\%)} & \textbf{AUC (\%)} \\
        \midrule
        r0 & AlexNet (scratch) & 78.54 & 62.06 & 74.41 & 67.67 & 85.53 \\
        r1 & AlexNet (pretrained) & 80.11 & 66.67 & 67.30 & 66.98 & 85.73 \\
        \rowcolor{bestgreen}
        \textbf{r2} & \textbf{EfficientNet-B3} & \textbf{79.97} & \textbf{63.71} & \textbf{78.20} & \textbf{70.21} & \textbf{88.93} \\
        \bottomrule
    \end{tabular}
    \caption{Lifecycle 1 test set metrics (binary classification, 699 images).}
\end{table}

\subsection{Confusion Matrices}

\begin{table}[H]
    \centering
    \renewcommand{\arraystretch}{1.2}
    \begin{tabular}{l|cc|cc|cc}
        \toprule
         & \multicolumn{2}{c|}{\textbf{r0 (AlexNet)}} & \multicolumn{2}{c|}{\textbf{r1 (Pretrained)}} & \multicolumn{2}{c}{\textbf{r2 (EfficientNet)}} \\
         & Pred: No & Pred: Yes & Pred: No & Pred: Yes & Pred: No & Pred: Yes \\
        \midrule
        Actual: No LS & 392 & 96 & --- & --- & 394 & 94 \\
        Actual: LS    & 54  & 157 & --- & --- & 46  & 165 \\
        \bottomrule
    \end{tabular}
    \caption{Confusion matrices for r0 and r2. r1 confusion matrix not recorded separately.}
\end{table}

% ============================================================================
\section{Analysis}
% ============================================================================

\subsection{What Worked}

\begin{enumerate}
    \item \textbf{Transfer learning significantly helps on small datasets:} AlexNet pretrained (r1) improved accuracy by +1.57\% over training from scratch, despite identical architecture.
    \item \textbf{EfficientNet-B3 dramatically improved recall:} Recall jumped from 74.41\% (r0) to 78.20\% (r2) --- 8 fewer missed landslides. For disaster monitoring, high recall is critical because a missed landslide can cost lives.
    \item \textbf{ROC-AUC improved substantially:} From 85.53\% to 88.93\% (+3.4\%), indicating the model's ranking ability is much better.
    \item \textbf{Dual-phase design works:} Phase 1 (CNN) and Phase 2 (HMM) integrate cleanly, producing actionable risk reports from raw satellite images.
\end{enumerate}

\subsection{Identified Weaknesses}

\begin{enumerate}
    \item \textbf{Overconfident predictions:} EfficientNet-B3 outputs softmax probabilities that are systematically overconfident --- predicting 95\% confidence on images where true accuracy is only $\sim$75\%. This makes confidence scores unreliable for operational decision-making.
    \item \textbf{Single-model limitations:} Each model has systematic blind spots. EfficientNet-B3 has high recall but lower precision (63.71\%); 94 false positives suggest the model is too aggressive. A single architecture cannot cover all failure modes.
    \item \textbf{Simplistic HMM:} With only 4--6 hidden states and 7 observation symbols (landslide type only), the HMM cannot distinguish between different trigger mechanisms (rainfall vs earthquake vs construction). This limits the quality of temporal forecasts.
    \item \textbf{Fixed 50\% threshold:} The default decision boundary is arbitrary and may not be optimal for the ensemble probability distribution.
\end{enumerate}

% ============================================================================
\section{Lifecycle 2 Preview}
% ============================================================================

Based on the weaknesses identified in Lifecycle 1, the next lifecycle will address three specific gaps:

\begin{enumerate}
    \item \textbf{Confidence Calibration (Temperature Scaling):}\\
          The overconfident softmax outputs from r2 make it impossible to set meaningful confidence thresholds for operational deployment (e.g., ``auto-alert if confidence $>$ 80\%''). Temperature scaling (Guo et al., ICML 2017) will be applied as a post-hoc calibration step to make confidence scores match true empirical accuracy.

    \item \textbf{Ensemble Methods:}\\
          Since EfficientNet-B3 and AlexNet make different types of errors, combining their predictions via weighted averaging should suppress individual model errors. We will also explore replacing AlexNet with a Vision Transformer (ViT-B/16) to create a CNN-Transformer hybrid ensemble that covers both local texture features (CNN) and global spatial relationships (Transformer).

    \item \textbf{Richer HMM Observations:}\\
          The current HMM observes only landslide type (7 symbols). By combining type with trigger mechanism into a single observation symbol (type $\times$ trigger = 42 symbols), the HMM can model finer-grained regimes such as ``monsoon-triggered debris flows'' vs ``earthquake-triggered rockfalls''.
\end{enumerate}

\noindent These improvements target calibration, model diversity, and temporal modeling --- the three weakest aspects of the current pipeline.

\end{document}
